\section{Residual Trust and Limitations}
\label{sec:limitations}

We collect the limitations stated elsewhere, ordered roughly by how much they matter to a user today.

\begin{enumerate}
  \item \textbf{No network.} Nothing travels over Tor. In the demonstration, buyer and vendor share one node and the letter is written to that node's own mailbox. Every network-privacy statement in this paper about Tor is a design statement.
  \item \textbf{Escrow is not two-party in any running node.} The node uses the merged single-process escrow, in which one process holds both secret keys. It is restricted to testnet4 in PR~\#2 and is labeled unsafe. The two-party protocol exists as a tested library in PR~\#3 with no transport.
  \item \textbf{The moderator is trusted to rule, and to rule in time.} A moderator that does not rule leaves a non-delivered buyer to lose at $h_C+1$. A moderator colluding with one party decides for that party. Equivocation is provable, not preventable.
  \item \textbf{The buyer must act before $h_D$.} The default outcome after $h_C$ is the vendor's, so a buyer who does not dispute non-delivery in time loses the payment.
  \item \textbf{Chain queries leak.} The Esplora server, and on testnet4 its fallback once used, sees every watched address and the node's IP address.
  \item \textbf{The mailbox key is not bound to the store.} Until the signed store record names the mailbox key, a buyer's letter is only as confidential as the lookup that told it which key to use.
  \item \textbf{At-rest encryption keeps its key beside the data.} It is not a defense against anyone who can read the data directory.
  \item \textbf{Local processes can act as the client.} The pairing token is available to any process on the machine.
  \item \textbf{No catalog freshness, completeness, replication, or spam cost.} A replica can serve stale listings or omit them; a store is available only from its own node today.
  \item \textbf{Payment rails beyond on-chain are not live.} BOLT 12 offers and BIP 352 codes are valid strings that nothing answers or scans; Payjoin is a library.
  \item \textbf{Amounts can identify listings, and payments can be linked.} Public prices let a chain observer guess what was bought from the payment amount, and because the wallet spends its own unconfirmed change, consecutive payments to different sellers can be chained together on chain.
  \item \textbf{Presigned fees.} Presigned escrow spends have fixed fees and no anchor output, so fee bumping is not possible and a fee spike near $h_C$ can let the vendor's claim race the buyer's dispute (finding M4).
  \item \textbf{No padding.} Letter sizes leak approximate content length.
  \item \textbf{Open review findings; no external review.} An internal adversarial review of the two-party escrow found the issues in Table~\ref{tab:review}; at \reviewcommit{} their status is as listed there. The code has not been audited by a third party, and the security arguments here are informal.
\end{enumerate}
